\chapter{Local Storage Declarations\label{chap:LocalStorageDeclarations}}

Local storage declarations introduce new variables into the \localstaticenvironmentterm.
A list of mutable (\Tvar) storage elements can be introduced without initializing any
of the storage elements.
ASL also supports tuples of (mutable or immutable) local storage declarations.

\ExampleDef{Local Storage Declarations}
\ASLListing{Local storage declarations}{LocalStorageDeclarations}{\definitiontests/LocalStorageDeclarations.asl}

\ChapterOutline
\begin{itemize}
  \item \FormalRelationsRef{Local Storage Declarations} defines the formal relations for local storage declarations;
  \item \SyntaxRef{Local Storage Declarations} defines the syntax of local storage declarations;
  \item \AbstractSyntaxRef{Local Storage Declarations} defines the abstract syntax of local storage declarations;
  \item \secref{Typing of Local Variable Declarations} defines the type rules of local variable declarations;
  \item \secref{Semantics of Local Variable Declarations} defines the dynamic semantics of local variable declarations;
  \item \secref{Typing of Tuple Declarations} defines the type rules of tuples of local storage declarations;
  \item \secref{Semantics of Tuple Declarations} defines the dynamic semantics of tuples of local storage declarations.
\end{itemize}

%%%%%%%%%%%%%%%%%%%%%%%%%%%%%%%%%%%%%%%%%%%%%%%%%%%%%%%%%%%%%%%%%%%%%%%%%%%%%%%%
\FormalRelationsDef{Local Storage Declarations}
\paragraph{Syntax:} Local storage declarations are grammatically derived from statements ($\Nstmt$)
  via dedicated rules (see \SyntaxRef{Local Storage Declarations}).

\paragraph{Abstract Syntax:} Local storage declarations are derived in abstract syntax by $\localdeclkeyword$ and $\localdeclitem$ and generated by $\buildlocaldeclkeyword$ and $\builddeclitem$
(see \AbstractSyntaxRef{Local Storage Declarations}).

\paragraph{Typing:}
\RenderRelation{annotate_local_decl_item}

This is used to annotate both single local storage declarations and tuples of local storage declarations.

\paragraph{Semantics:}
\RenderRelation{eval_local_decl}
That is, the right-hand side of the declaration
has already been evaluated, yielding $\vm$ (see, for example, \SemanticsRuleRef{SDecl}).

This is used to evaluate both single local storage declarations and tuples of local storage declarations.

%%%%%%%%%%%%%%%%%%%%%%%%%%%%%%%%%%%%%%%%%%%%%%%%%%%%%%%%%%%%%%%%%%%%%%%%%%%%%%%%
\SyntaxDef{Local Storage Declarations}
\hypertarget{def-localdeclarationkeyword}{}
\hypertarget{def-localdeclarationitem}{}
A \localdeclarationkeyword\ is either \texttt{var} or \texttt{let}.
A \localdeclarationitem\ is an element derived from $\Ndeclitem$.
\hypertarget{def-localdeclaration}{}
A \localdeclaration\ consists of a \localdeclarationitem\ and a \localdeclarationkeyword.

Declaring a local storage element is done via the following grammar rules:
\begin{flalign*}
\Nstmt \derives \ & \Nlocaldeclkeyword \parsesep \Ndeclitem \parsesep \option{\Nasty} \parsesep \Teq \parsesep \Nexpr \parsesep \Tsemicolon &\\
|\ & \Tvar \parsesep \Ndeclitem \parsesep \Nasty \parsesep \Tsemicolon &\\
|\ & \Tvar \parsesep \Clisttwo{\Tidentifier} \parsesep \Nasty \parsesep \Tsemicolon &\\
\end{flalign*}

\begin{flalign*}
\Nlocaldeclkeyword \derives \ & \Tlet \;|\; \Tvar &\\
\Ndeclitem \derives\
   & \Tidentifier &\\
|\ & \Plisttwo{\Nignoredoridentifier} &\\
\Nignoredoridentifier \derives \ & \Tminus \;|\; \Tidentifier &
\end{flalign*}

%%%%%%%%%%%%%%%%%%%%%%%%%%%%%%%%%%%%%%%%%%%%%%%%%%%%%%%%%%%%%%%%%%%%%%%%%%%%%%%%
\AbstractSyntaxDef{Local Storage Declarations}
\RenderTypes[remove_hypertargets]{local_decl_keyword_and_item}


\RequirementDef{DiscardingLocalStorageDeclarations}
Local storage declarations must bind at least one name.
All the local storage declarations in \listingref{local-storage-discards} are illegal, as they discard all declared storage elements.
\ASLListing{Illegal local storage declarations that discard all storage elements}{local-storage-discards}{\syntaxtests/GuideRule.DiscardingLocalStorageDeclarations.asl}

\ASTRuleDef{LocalDeclKeyword}
\RenderRelation{build_local_decl_keyword}
\RenderProseAndFormally{build_local_decl_keyword}

\ASTRuleDef{DeclItem}
\RenderRelation{build_decl_item}
\RenderProseAndFormally{build_decl_item}

\ASTRuleDef{IgnoredOrIdentifier}
\RenderRelation{build_ignored_or_identifier}
\RenderProseAndFormally{build_ignored_or_identifier}

\section{Typing of Local Variable Declarations\label{sec:Typing of Local Variable Declarations}}
\TypingRuleDef{LDVar}
\ExampleDef{Well-typed Local Variable Declarations}
In \listingref{ldvar}, the statement \texttt{let x = 3;} is legal, since
\texttt{x} is not defined elsewhere. It is added to the type environment
with the type inferred type \texttt{integer{3}}.
\ASLListing{A local storage declaration}{ldvar}{\typingtests/TypingRule.LDVar.asl}

\RenderProseAndFormally{annotate_local_decl_item_var}
\CodeSubsection{\LDVarBegin}{\LDVarEnd}{../Typing.ml}

\TypingRuleDef{CheckIsNotCollection}%
\RenderRelation{check_is_not_collection}

\ExampleDef{Check that a Type is not a Collection}
In \listingref{checkisnotcollection}, the statement
\verb|var test: MyCollection;| fails with a \typingerrorterm{} because
\TypingRuleRef{LDVar} calls \TypingRuleRef{CheckIsNotCollection}.
\ASLListing{Declaring a local variable with a collection type}{checkisnotcollection}{\typingtests/TypingRule.CheckIsNotCollection.asl}


\RenderProseAndFormally{check_is_not_collection}

\CodeSubsection{\CheckIsNotCollectionBegin}{\CheckIsNotCollectionEnd}{../Typing.ml}

\section{Semantics of Local Variable Declarations\label{sec:Semantics of Local Variable Declarations}}
\SemanticsRuleDef{LDVar}
\ExampleDef{Evaluation of a Local Variable Declaration}
The statement \texttt{var x = 3;} in \listingref{ldvar} binds \texttt{x}
to the evaluation of \texttt{3} in $\env$.


\RenderProseAndFormally{eval_local_decl_LDI_Var}
\CodeSubsection{\EvalLDVarBegin}{\EvalLDVarEnd}{../Interpreter.ml}

\section{Typing of Tuple Declarations\label{sec:Typing of Tuple Declarations}}
\TypingRuleDef{LDTuple}
\ExampleDef{Well-typed Tuple Declarations}
\ASLListing{Declaring a tuple in the local storage}{typing-ldtuple}{\typingtests/TypingRule.LDTuple.asl}

\RenderProseAndFormally{annotate_local_decl_item_tuple}
\CodeSubsection{\LDTupleBegin}{\LDTupleEnd}{../Typing.ml}

\TypingRuleDef{AddLocalVars}
\RenderRelation{add_local_vars}

See \ExampleRef{Well-typed Tuple Declarations}.

\RenderProseAndFormally{add_local_vars}

\section{Semantics of Tuple Declarations\label{sec:Semantics of Tuple Declarations}}
\SemanticsRuleDef{LDTuple}
\ExampleDef{Evaluation of Tuple Declarations}
In \listingref{semantics-ldtuple},
\texttt{var (x,y,z) = (1,2,3);} binds \texttt{x} to the evaluation of
\texttt{1}, \texttt{y} to the evaluation of \texttt{2}, and \texttt{z} to the
evaluation of \texttt{3} in $\env$.
\ASLListing{Evaluating a tuple declaration in the local storage}{semantics-ldtuple}{\semanticstests/SemanticsRule.LDTuple.asl}


\RenderProseAndFormally{eval_local_decl_LDI_Tuple}
\CodeSubsection{\EvalLDTupleBegin}{\EvalLDTupleEnd}{../Interpreter.ml}

\SemanticsRuleDef{DeclareLDITuple}
\RenderRelation{declare_ldi_tuple}

See \ExampleRef{Evaluation of Tuple Declarations}.


\RenderProseAndFormally{declare_ldi_tuple}
